\chapter{Procedural Universes}

\chapterepigraph{%
  A procedurally generated world is not a copy of the real world.\\
  It is a theory of the real world:\\
  a set of rules that produces worlds\\ like this one.
}{Flyxion, \textit{Semantic Infrastructure}}

\sidenote{Connects to\\RSVP field\\equations (Ch.\,18)\\and constraint\\ecology (Ch.\,26)}

Procedural generation — the creation of content through algorithms
rather than explicit design — is one of the most intellectually rich
areas of game design and computer graphics. This chapter examines
it through the lens of the accessibility framework.

\section{Generation}

A procedural generator is a constraint satisfaction system: given a
set of high-level constraints (``generate a dungeon that is winnable,
interesting, and of approximately level 5 difficulty''), it produces
specific instantiations (a particular dungeon layout) that satisfy them.

\begin{definition}{Procedural Generator}{proc-gen}
  A \emph{procedural generator} is a map $G : \Theta \times \mathcal{R}
  \to W$ where $\Theta$ is a parameter space (constraints), $\mathcal{R}$
  is a randomness source (seeds), and $W$ is a space of worlds.
  The generator is \emph{good} if: (1) all generated worlds satisfy
  the constraints, and (2) the distribution over worlds is diverse
  (high entropy given the constraints).
\end{definition}

\section{Emergence}

The most interesting procedural worlds exhibit emergence: global
structure that arises from local rules without being explicitly
specified. Cellular automata (Conway's Game of Life), L-systems
(plant growth), and agent-based simulations all produce emergent
structures.

Emergence is the procedural version of the RSVP accessibility
relaxation: local constraint satisfaction rules (cells reproduce
if they have 2--3 neighbors) produce global accessible patterns
(gliders, oscillators, stable configurations) that persist because
they are attractors in the accessibility landscape of the rule system.

\section{Constraint Satisfaction in World Construction}

The constraint satisfaction approach to procedural generation makes
the accessibility framework explicit:

\begin{itemize}
  \item \textbf{High-level constraints:} the world must be habitable,
        have a balanced ecosystem, have navigable terrain.
  \item \textbf{Low-level generation:} tile placement, creature
        population, resource distribution.
  \item \textbf{Validation:} constraint violation check — does the
        generated world satisfy the high-level constraints?
  \item \textbf{Repair:} if not, modify the low-level generation
        to reduce constraint violation (constraint descent).
\end{itemize}

This is the Spherepop evaluation of world construction: the high-level
constraints are the outer spheres; the low-level details are the inner
spheres; generation is evaluation (collapse) from inside out.

\begin{exercise}
  Describe the constraint space for generating a believable ecological
  system: predator-prey relationships, carrying capacity, seasonal
  variation, evolutionary pressure. What are the minimum constraints
  needed to produce ecologically plausible behavior? What emergent
  structures arise from these constraints?
\end{exercise}

\begin{exercise}
  The No Man's Sky universe (18 quintillion planets) is procedurally
  generated from a seed plus rules. Using the accessibility framework,
  explain why many players found the universe ``empty'' despite its size.
  What constraints were missing? What accessibility landscape would
  a more engaging procedural universe require?
\end{exercise}
